\section{Proposed Solution}
\label{sec:proposed}

An analyst who reviews a BRSR report must connect what the company writes with the numbers the company filed. GRASP supports this work by organising the review around one shared source, the structured BRSR report of each company. Text segments and facts are stored together with the company and reporting details. A \emph{text segment} is the selected portion of a report's text that is analysed, and a \emph{fact} is one reported value with the concept, unit, and period. Because both are stored under the same report, an analyst can move from a flagged text segment or a score component to the underlying disclosure. The central idea is to connect language analysis with reported evidence and to present all results in one review workflow.

The design responds to the three questions raised in Section~\ref{sec:introduction}. For the first question, evidence links connect each text segment with the related facts. For the second question, the source of every training label is recorded, so the way each dataset row received a principle label can be checked. For the third question, a visible score breakdown and a sector comparison explain the position of a company. Each model, rule, and calculation has a defined role in this design.

The workflow is built from the seven modules shown in Figure~\ref{fig:architecture}. The first three modules prepare the data and train the classifier. Report Acquisition and Structured Storage creates the common data record. Text Segment Selection and Weak-Label Construction removes unusable text and assigns a \emph{weak label} to each remaining text segment. A weak label is a principle label taken automatically from the principle tags and concept mappings of the report, including the documented AI-assisted assignments. Domain Adaptation and Principle Classification then prepares BERT for BRSR language and trains a classifier for the nine reporting principles.

The next three modules analyse the stored text segments and facts for different purposes. Vagueness Detection and Report-Consistency Checks has two parts. Word lists and sentence similarity identify vague wording, while numerical and assurance checks create separate review records that contain the relevant claim and the reported value. Principle Mapping and the Zero-Shot Comparator uses the probabilities from the classifier to summarise how strongly each principle is supported, and a zero-shot method serves as the comparison for evaluation. Transparency Score and Peer Comparison combines four measures, namely Completeness, Specificity, Clarity, and SDG breadth, and then compares each company with other companies in the same sector. The principle probabilities and the numerical review flags are not part of these four measures. These probabilities and flags remain separate outputs.

The last module, Service Layer and Runtime Traceability, makes the outputs available through an application programming interface (API). These outputs are company records, review flags, principle summaries, score breakdowns, and peer comparisons. Results can be read at four levels: one company, selected companies within a sector, the complete financial-sector group, and all recorded sectors. Where a company has two reports from different years, the results can also be compared year by year. Section~\ref{sec:methodology} explains how each module works, and Section~\ref{sec:evaluation} defines the evidence used to assess the outputs.
